
% =============================================================================
% Table 01 — Forms of Invariance
% Book I — Ontology of Viability
% Chapter 02 — Invariance
%
% Purpose:
%   To provide an analytical classification of possible forms of
%   invariance in relation to different kinds of determinations,
%   relations, and transformations.
%
% Conceptual Policy:
%   1. The categories are analytical distinctions, not necessarily
%      mutually exclusive ontological kinds.
%   2. No form of invariance is presumed to apply to every existent.
%   3. Invariance is distinct from identity, persistence, continuity,
%      stability, and immutability.
%   4. Each form requires specification of its domain and respect.
%   5. The classification remains open to refinement.
% =============================================================================

\begin{table}[htbp]
    \centering
    \small
    \renewcommand{\arraystretch}{1.35}
    \setlength{\tabcolsep}{5pt}

    \caption{Analytical Forms of Invariance}
    \label{tab:forms-of-invariance}

    \begin{tabular}{
        p{2.7cm}
        p{3.1cm}
        p{3.4cm}
        p{4.0cm}
    }
        \hline
        \textbf{Form of Invariance}
        &
        \textbf{What Remains Invariant}
        &
        \textbf{Relevant Transformation}
        &
        \textbf{Ontological Qualification}
        \\
        \hline

        Determinational Invariance
        &
        A particular determination or property
        &
        Transformations affecting other determinations
        &
        A determination may remain preserved even when other aspects
        of the existent undergo change.
        \\
        \hline

        Structural Invariance
        &
        An organizational pattern or structural relation
        &
        Transformations affecting components or their configurations
        &
        Structural preservation does not necessarily imply preservation
        of the components themselves.
        \\
        \hline

        Relational Invariance
        &
        A specified relation among relata
        &
        Transformations affecting the relata or their configurations
        &
        A relation may remain preserved despite changes in the
        determinations of its relata.
        \\
        \hline

        Compositional Invariance
        &
        A specified compositional organization
        &
        Transformations affecting components, their arrangement,
        or their interactions
        &
        The persistence of a compositional pattern must be distinguished
        from the persistence of the same material components.
        \\
        \hline

        Quantitative Invariance
        &
        A specified quantitative magnitude or measure
        &
        Transformations that may alter the configuration of the existent
        &
        Quantitative preservation does not establish qualitative,
        structural, or numerical identity.
        \\
        \hline

        Qualitative Invariance
        &
        A specified qualitative determination
        &
        Transformations that may alter other qualitative or quantitative
        aspects
        &
        The preservation of a quality must be distinguished from the
        preservation of the entire configuration.
        \\
        \hline

        Configurational Invariance
        &
        A specified configuration or pattern of determinations
        &
        Transformations that may alter the underlying elements or
        their relations
        &
        The relevant criterion of configuration must be explicitly
        specified; similarity alone is insufficient.
        \\
        \hline

        Relational-Structural Invariance
        &
        A specified organization of relations
        &
        Transformations that preserve relational organization while
        modifying particular relata or determinations
        &
        This form concerns the preservation of relational organization,
        not necessarily the preservation of every individual relation.
        \\
        \hline

    \end{tabular}

    \vspace{0.5em}

    \begin{minipage}{0.95\textwidth}
        \footnotesize
        \textit{Note:}
        The forms listed above are analytical categories intended
        to distinguish different possible respects of invariance.
        They are not asserted to be mutually exclusive, exhaustive,
        or ontologically independent.
        A single transformation may preserve more than one kind
        of determination or relation, while altering others.
        The classification does not establish that any particular
        invariant exists in every existent.
    \end{minipage}

\end{table}
